\documentclass[addpoints]{exam}
% \documentclass[addpoints,12pt]{exam}

\usepackage[slovene]{babel}
\usepackage{amsfonts}
\usepackage{amssymb}
\usepackage{amsmath}
\usepackage{float}
\usepackage{graphicx}
\usepackage{enumitem}
\usepackage{color}
\usepackage{eurosym}


%Komentar naslednje vrstice glede na verzijo testa -- rešitve ali prostor za odgovore:
\printanswers

% \linespread{2}


\pointsinrightmargin
\marginbonuspointname{}


\renewcommand{\solutiontitle}{\noindent\textbf{Rešitve in točkovanje:}\par\noindent}

\colorgrids
\definecolor{GridColor}{gray}{0.7}
\setlength{\gridsize}{7mm}

\extrawidth{5mm}
\setlength{\rightpointsmargin}{1.5cm}

\begin{document}

\runningfooter{}{}{}

\begin{center}
    \fbox{\fbox{\parbox{15cm}{\centering
    Tretje pisno ocenjevanje znanja -- ponovno \\ 1. letnik -- test E \\ 17. marec 2025 \\ (Deljivost; Racionalna števila)}}}
\end{center}

\vspace{0.1in}
\makebox[\textwidth]{Ime in priimek, oddelek:\enspace\hrulefill}


\begin{center}
    \chqword{Naloga}
    \chpword{Možne točke}
    \chbpword{Dodatne točke}
    \chsword{Dosežene točke}
    \chtword{Skupaj}  
    \combinedpointtable[h][questions]
    % \combinedgradetable[h][questions]

\end{center}

\vspace{0.1in}


\begin{questions}

    % 1. vprašanje

    \question[4]
    Določite največji skupni delitelj in najmanjši skupni večkratnik izrazov $x^3-4x$ in $x^3-8$.

    \begin{solution}[6cm]
        Za pravilno faktorizacijo izrazov $x(x-2)(x+2)$ in $(x-2)(x^2+2x+4)$ po $1$ točka. 
        Za zapis največjega skupnega delietlja $x-2$ $1$ točka in najmanjšega skupnega večkratnika $x(x-2)(x+2)(x^2+2x+4)$ $1$ točka.
    \end{solution}
    
    
    % 2. vprašanje
    \question
    Zapišite po besedilu in izračunajte število $x$ oziroma $y$.
    \begin{parts}
    \part[2]
        Če število $x$ delimo s številom $123$, dobimo količnik $7$ in ostanek $54$.
    
    \begin{solution}[4cm]
        Za pravilen zapis $x=123\cdot 7+53$ $1$ točka in izračun $x=914$ $1$ točka.
    \end{solution}

    \part[3]
        Če število $5613$ delimo s številom $y$, dobimo količnik $17$ in ostanek $292$.
    
    \begin{solution}[3cm]
        Za pravilen zapis $5613=y\cdot 17+292$ $1$ točka; za izračun $y=313$ $2$ točki.
    \end{solution}

        
    \end{parts}




     \newpage

    % 3. vprašanje
    % \question[4]
    % Produkt dveh števil je $192$, njun največji skupni delitelj pa $4$. Poiščite števili.
    
    % \begin{solution}[7cm]
    %     \begin{align*}
    %         1-(1+x^8)(1+x^4)(1+x^2)(1+x)(1-x)&=1-(1+x^8)(1+x^4)(1+x^2)(1-x^2)= \\
    %             &=1-(1+x^8)(1+x^4)(1-x^4)= \\
    %             &=1-(1+x^8)(1-x^8)= \\
    %             &=1-(1-x^{16})= \\
    %             &=1-1+x^{16}= \\
    %             &=x^{16}
    %     \end{align*}
    %     Za pravilno uporabo pravila razlike kvadratov $1$ točka, za pravilno upoštevan negativen predznak $1$ točka, končen rezultat $1$ točka.
    % \end{solution}


    % \begin{description}
    %     \item[a] $(x-a)^2+(y-b)^2-c^2=0$ 
    %     \item[b] $d^2(x-a)^2+f^2(y-b)^2-(df)^2=0$
    %     \item[c] $S(a,b)$
    %     \item[d] $e^2=d^2-f^2;\ d>f$
    % \end{description}
    % Krožnica: \fillin[a c][2cm]
    % \\ Elipsa:  \fillin[b c d][2cm]


        % 4. vprašanje
        \question[3]
        Ali je število $559$ praštevilo? Utemeljite.
        
        \begin{solution}[7cm]
            Za zapis pravila za preverjanje oziroma njegova uporaba $1$ točka; za preveritev deljivosti z vsemi potrebnimi praštevili $1$ točka;
            za ustrezen zaključek, da $559$ ni praštevilo, $1$ točka.
        \end{solution}
    

    % \newpage
    % 5. vprašanje
    \question
    Izračunajte vrednost izraza.
    \begin{parts}
        \part[5]
        $\dfrac{\frac{11}{6}+3\frac{3}{8}}{\frac{19}{12}}-2\frac{1}{12}\cdot \frac{18}{5}$

        \begin{solution}[7cm]
            \begin{align*} 
                \dfrac{\frac{11}{6}+3\frac{3}{8}}{\frac{19}{12}}-2\frac{1}{12}\cdot \frac{18}{5} &=\dfrac{\frac{11}{6}+\frac{27}{8}}{\frac{19}{12}}-\frac{25}{12}\cdot\frac{18}{5}= \\
                                                                                            &=\dfrac{\frac{44}{24}+\frac{81}{24}}{\frac{19}{12}}-\frac{5}{2}\cdot 3= \\
                                                                                            &=\dfrac{\frac{125}{24}}{\frac{19}{12}}-\frac{15}{2}= \\
                                                                                            &=\frac{125\cdot 12}{24\cdot 19}-\frac{15}{2}= \\
                                                                                            &=\frac{125}{38}-\frac{285}{38}= \\
                                                                                            &=\frac{160}{38}= \\
                                                                                            &=\frac{80}{19}
            \end{align*}
            Za pravilno množenje ulomkov $1$ točka, pravilno seštevanje ulomkov $1$ točka, pravilno razreševanje dvojnih ulomkov $1$ točka, pravilno odštevanje ulomkov $1$ točka, pravilen okrajšan rezultat $1$ točka.
        \end{solution}


        \part[5]
        $5.\overline{6}:(1.5-1.\overline{3}) - 0.3\overline{72}\cdot 110 $

        \begin{solution}[4cm]
            \begin{align*}
                5.\overline{6}:(1.5-1.\overline{3}) - 0.3\overline{72}\cdot 110&=5\frac{2}{3}:(1\frac{1}{2}-1\frac{1}{3})- \frac{41}{110}\cdot 110 = \\
                                                                            &=\frac{17}{3}:\frac{1}{6}-41= \\
                                                                            &=\frac{17}{3}\cdot \frac{6}{1}-41= \\
                                                                            &=34-41= \\
                                                                            &= -7
            \end{align*}
            Za pravilno zapisano število $0.3\overline{72}$ z ulomkom $1$ točka, pravilno zapisani števili $5.\overline{6}$ in $1.\overline{3}$ z ulomkoma $1$ točka, pravilno množenje in deljenje ulomkov $1$ točka, pravilno odštevanje ulomkov $1$ točka, končen rezultat $1$ točka.
        \end{solution}

    \end{parts}



    \newpage
    % 6. vprašanje

    \question
    Okrajšajte ulomke.
    \begin{parts}
        

        \part[4]
        $\dfrac{x+2}{(x-2)^{-1}}:\left(\dfrac{1}{x^2-4}\right)^{-1}$

        \begin{solution}[7cm]
            \begin{align*}
                \dfrac{x+2}{(x-2)^{-1}}:\left(\dfrac{1}{x^2-4}\right)^{-1}&= \dfrac{(x-2)(x+2)}{1}\cdot\dfrac{1}{x^2-4}= \\
                                                                            &= \dfrac{x^2-4}{x^2-4}= \\
                                                                            &= 1
            \end{align*}
            Za pravilno upoštavanje obratne vrednosti $1$ točka; pravilno deljenje ulomkov $1$ točka; pravilna faktorizacija oziroma razčlenjevanje $1$ točka; pravilen okrajšan rezultat $1$ točka.
        \end{solution}


        \part[5]
        $\dfrac{7^{3n-3}+3\cdot 7^{3n-2}-7^{3n-4}}{7^{3n-2}-7^{3n-1}}$

        \begin{solution}[7cm]
            \begin{align*}
                \dfrac{7^{3n-3}+3\cdot 7^{3n-2}-7^{3n-4}}{7^{3n-2}-7^{3n-1}} &= \dfrac{7^{3n-4}\left(7^1+3\cdot 7^2-1\right)}{7^{3n-2}\left(1-7^1\right)}= \\
                                                                            &= \dfrac{7^{3n-4}\left(7+3\cdot 49-1\right)}{7^{3n-2}\left(1-7\right)}= \\
                                                                            &= \dfrac{7^{3n-4}\cdot 153}{7^{3n-2}\cdot(-6)}= \\
                                                                            &= \dfrac{7^{3n-4}}{7^{3n-2}}\cdot\dfrac{-51}{2}= \\
                                                                            &= -\dfrac{51}{2\cdot 7^2}
            \end{align*}
            Za pravilno izpostavljanje skupnega faktorja v števcu in imenovalcu po $1$ točka, pravilno seštevanje $1$ točka, pravilno krajšanje potenc $1$ točka, zapis pravilno okrajšanega ulomka $1$ točka.
        \end{solution}


        \part[5]
        $\dfrac{(\frac{1}{3})^{-1}a^2\cdot 10^{-3}(a^{-2})^4 (10^{-4})^3 (a^0+11b^0)}{(6a^{-3}10^{-8})^2}; \quad a,b\neq 0$

        \begin{solution}[3cm]
            \begin{align*}
                \dfrac{(\frac{1}{3})^{-1}a^2\cdot 10^{-3}(a^{-2})^4 (10^{-4})^3 (a^0+11b^0)}{(6a^{-3}10^{-8})^2} &= \dfrac{3a^2\cdot 10^{-3}a^{-8} 10^{-12} (1+11\cdot 1)}{6^2a^{-6}10^{-16}}= \\
                                                                                                                &= \dfrac{3\cdot 12 a^{2-8}10^{-3-12}}{36a^{-6}10^{-16}}= \\
                                                                                                                &= a^{-6-(-6)}10^{-15-(-16)}= \\
                                                                                                                &= a^0 10= \\
                                                                                                                &= 10
            \end{align*}
            Za uporabo pravila potenciranja potence $1$ točka, pravilno obratna vrednost ulomka $1$ točka, uporaba pravila množenja potenc z enakimi osnovami $1$ točka, pravilno upoštevanje $x^0=1$ $1$ točka, končen rezultat $1$ točka.
        \end{solution}


    \end{parts}



    \newpage

        % 7. vprašanje

        \question
        Poenostavite izraze.
        \begin{parts}
            

    
            \part[5]
            $\dfrac{2x^{-2}}{1-2x^{-1}+x^{-2}}+\dfrac{x^{-1}}{1-x^{-1}}; \quad x\neq 1$
    
            \begin{solution}[7cm]
                \begin{align*}
                    \dfrac{2x^{-2}}{1-2x^{-1}+x^{-2}}+\dfrac{x^{-1}}{1-x^{-1}} &= \dfrac{\frac{2}{x^2}}{1-\frac{2}{x}+\frac{1}{x^2}}+\dfrac{\frac{1}{x}}{1-\frac{1}{x}}= \\
                                                                            &= \dfrac{\frac{2}{x^2}}{\frac{x^2-2x+1}{x^2}}+\dfrac{\frac{1}{x}}{\frac{x-1}{x}}= \\
                                                                            &= \dfrac{2x^2}{(x^2-2x+1)x^2}+\dfrac{x}{(x-1)x}= \\
                                                                            &= \dfrac{2}{(x-1)^2}+\dfrac{1}{x-1}= \\
                                                                            &= \dfrac{2+(x-1)}{(x-1)^2}= \\
                                                                            &= \dfrac{x+1}{(x-1)^2}
                \end{align*}
                Za pravilen zapis obratnih vrednosti $1$ točka, pravilno seštevanje in odštevanje ulomkov $1$ točka, pravilno razreševanje dvojnih ulomkov $1$ točka, pravilno krajšanje ulomkov $1$ točka, končen rezultat $1$ točka.
            \end{solution}
    
    
            \part[5]
            $\dfrac{(a-1)^{-1}-(a+1)^{-1}}{a+\frac{a}{a^2-1}}; \quad a\neq 0$
    
            \begin{solution}[7cm]
                \begin{align*}
                    \dfrac{(a-1)^{-1}-(a+1)^{-1}}{a+\frac{a}{a^2-1}} &= \dfrac{\frac{1}{a-1}-\frac{1}{a+1}}{\frac{a(a^2-1)+a}{a^2-1}}= \\
                                                                    &= \dfrac{\frac{a+1-(a-1)}{a^2-1}}{\frac{a^3-a+a}{a^2-1}}= \\
                                                                    &= \dfrac{\frac{2}{a^2-1}}{\frac{a^3}{a^2-1}}= \\
                                                                    &= \dfrac{2(a^2-1)}{a^3(a^2-1)}= \\
                                                                    &= \dfrac{2}{a^3}
                \end{align*}
                Za pravilen zapis obratnih vrednosti ulomkov $1$ točka, pravilno seštevanje in odštevanje ulomkov $1$ točka, pravilno razrešen dvojni ulomek $1$ točka, pravilno krajšanje ulomkov $1$ točka, končen rezultat $1$ točka.
            \end{solution}
    

            % \part[4]
            % $\dfrac{x+4}{x+5}-\dfrac{2x-10}{x^2-x-12}:\dfrac{x^2-25}{x-4}; \quad x\neq -3$
    
            % \begin{solution}[3cm]
            %     \begin{align*}
            %         \dfrac{x+4}{x+5}-\dfrac{2x-10}{x^2-x-12}:\dfrac{x^2-25}{x-4} &= \dfrac{x+4}{x+5}-\dfrac{2(x-5)}{(x-4)(x+3)}\cdot\dfrac{x-4}{(x-5)(x+5)}= \\
            %                                                                     &= \dfrac{x+4}{x+5}-\dfrac{2}{(x+3)(x-5)}= \\
            %                                                                     &= \dfrac{(x+4)(x+3)-2}{(x+5)(x+3)}= \\
            %                                                                     &= \dfrac{x^2+7x+10}{(x+5)(x+3)}= \\
            %                                                                     &= \dfrac{(x+5)(x+2)}{(x+5)(x+3)}= \\
            %                                                                     &=\dfrac{x+2}{x+3}
            %     \end{align*}
            %     Za pravilno faktorizacijo izrazov $1$ točka, pravilno deljenje ulomkov $1$ točka, pravilno odštevanje ulomkov $1$ točka, pravilno zapisan okrajšan rezultat $1$ točka.
            % \end{solution}
    

        \end{parts}


    %  \newpage






     \newpage

     \noindent Ime in priimek: \underline{~~~~~~~~~~~~~~~~~~~~~~~~~~~~~~~~~~~~~~~~~~~~~~~~}

     ~\\

    %dodatno vprašanje
    \bonusquestion[2]
    \textit{Dodatna naloga:} Zapišite število $524$ v osmiškem številskem sestavu.

    \begin{solution}[7cm]
        \begin{align*}
            524 &= 8\cdot 65 + 4 \\
            65 &= 8\cdot 8 + 1 \\
            8 &= 8\cdot 1 + 0 \\
            1 &= 8\cdot 0 + 1
        \end{align*}
        Za pravilno zapisan algoritem $1$ točka, zapis števila v osmiškem sistemu $524_{(10)}=1014_{(8)}$ $1$ točka.
    \end{solution}




    % \newpage
        

    % dodatno vprašanje
    \bonusquestion[3]
    \textit{Dodatna naloga:} Okrajšajte ulomek: $\dfrac{b^{2x-y}+b^x+b^{x-y+2}}{b^{2x+y}+b^{x+2y}+b^{x+y+2}}$.

    \begin{solution}[7cm]
        \begin{align*}
            \dfrac{b^{2x+y}+b^{x+2y}+b^{x+y+2}}{b^{2x-y}+b^x+b^{x-y+2}} &= \dfrac{b^{x-y}(b^x+b^y+b^2)}{b^{x+y}(b^x+b^y+b^2)}= \\
                                                                        &= \dfrac{b^{x-y}}{b^{x+y}}= \\
                                                                        &= b^{-2y}
        \end{align*}
        Za pravilno izpostavljanje skupnega faktorja $1$ točka; za pravilno krajšanje $1$ točka; končen rezultat $1$ točka.
    \end{solution}




\end{questions}



\end{document}